\chapter{Methods for Regression Evaluation}
\label{chap:regression-methods}

\section{Regression Framework}
\label{sec:reg-framework}

\subsection{Linear Regression}
The linear regression model is the foundation for analyzing the relationship between a continuous outcome $Y$ and a set of predictors $\mathbf{X} = \{X_1, X_2, \dots, X_p\}$ \cite{mccullagh1989}. The model is defined as:
\begin{equation}
    Y = \beta_0 + \sum_{j=1}^{p} \beta_j X_j + \epsilon, \quad \epsilon \sim \mathcal{N}(0, \sigma^2)
\end{equation}
where $\beta_j$ are the regression coefficients and $\epsilon$ represents the error term.

\subsection{Logistic Regression}
For binary outcomes $Y \in \{0, 1\}$, we employ the logistic regression model \cite{hosmer2013}:
\begin{equation}
    \text{logit}(P(Y=1|\mathbf{X})) = \ln\left(\frac{P(Y=1|\mathbf{X})}{1-P(Y=1|\mathbf{X})}\right) = \beta_0 + \sum_{j=1}^{p} \beta_j X_j
\end{equation}

\subsection{Inferential Quantities of Interest}
Our evaluation focuses on the preservation of the estimated coefficients $\hat{\beta}_j$ and their associated 95\% confidence intervals (CIs). If synthetic data is to be used as a proxy for real data, the synthetic coefficients $\hat{\beta}_j^{(S)}$ should be unbiased estimates of the original coefficients $\hat{\beta}_j^{(R)}$, and their CIs should overlap substantially.

\section{Data Framework}
\label{sec:reg-data-framework}

\subsection{Original Datasets and Synthetic Datasets}
We generate Original Datasets (OD) through a controlled simulation process. From each OD, we create one corresponding Synthetic Dataset (SD) using a specified generator.

\subsection{General Comparison Strategy: OD versus SD}
The utility of the SD is measured by comparing the results of the same regression model fitted to both the OD and the SD. This "one-to-one" comparison ensures that any observed degradation is attributable solely to the synthesis process.

\subsection{Role of Repeated Simulation Scenarios}
To ensure statistical robustness, each unique scenario (combination of factors) is repeated $M = 1,000$ times. This allows us to calculate the average performance and the variability of the evaluation metrics across independent runs.

\section{Data Generation Mechanisms}
\label{sec:reg-generation-mechanisms}

\subsection{Simulation Factors}
The simulation space is defined by four main factors: sample size ($N$), number of predictors ($p$), correlation between predictors ($\rho$), and error variability ($\sigma^2$ for linear) or outcome prevalence (for logistic).

\subsection{Sample Size, Number of Predictors, and Variable Type}
We test three levels of sample size ($N \in \{100, 500, 1000\}$) and three levels of dimensionality ($p \in \{2, 5, 10\}$). The variable types are divided into two blocks: all-continuous predictors and all-binary predictors.

\subsection{Correlation Structure and Error Variability}
Predictors are generated from a multivariate normal distribution (then discretized for binary scenarios) with a covariance matrix $\Sigma$ where all off-diagonal entries are set to $\rho \in \{0.0, 0.3, 0.6\}$. For linear models, the error variance $\sigma^2$ is varied in $\{0.5, 1.0, 2.0\}$.

\subsection{Regression Coefficients and Outcome Generation}
The coefficients are fixed at $\beta_j = (-1)^j$ to ensure a mix of positive and negative relationships. The outcome $Y$ is then generated based on the linear predictor $\mathbf{X}\beta$ plus the stochastic component.

\section{Synthetic Data Generators}
\label{sec:reg-generators}

\subsection{Generators for Predictors}
We evaluate two primary approaches for synthesizing the predictors $\mathbf{X}$:
\begin{itemize}
    \item \textbf{CART:} Non-parametric classification and regression trees.
    \item \textbf{Parametric:} Normal distribution for continuous variables and logistic regression for binary ones.
\end{itemize}

\subsection{Generators for Outcomes}
The outcome $Y$ is synthesized last, conditioned on all previously generated predictors $\mathbf{X}$, using the same family of methods (CART or Parametric).

\subsection{Combined Generation Arms}
The study evaluates "homogeneous" arms where the same method family is used for both predictors and outcomes (e.g., CART-X + CART-Y vs. Norm-X + Norm-Y).

\subsection{Conceptual Differences between the Generators}
Parametric generators assume the underlying data follows a specific distribution, whereas CART generators are model-agnostic and can capture non-linearities and interactions, albeit at the risk of overfitting or introducing "blocky" artifacts.

\section{Experimental Design}
\label{sec:reg-experimental-design}

\subsection{Scenario Grid}
The full experimental design comprises 1,800+ scenarios, resulting from the Cartesian product of all simulation factors.

\subsection{Number of Repetitions and Successful Fits}
Each scenario is executed until $M=1,000$ successful fits are obtained. Scenarios that fail to converge (common in logistic regression with small $N$ and high $p$) are flagged but excluded from the final metric calculation to ensure comparability.

\subsection{Computational Workflow}
The workflow is implemented in R (using \texttt{synthpop}) for data generation and synthesis, and Python for the evaluation and visualization of results.

\subsection{Quality-Control and Fit-Validation Rules}
We apply strict quality control to ensure that the original data is "well-behaved" (no perfect collinearity) before synthesis, ensuring that any failures in the SD are truly reflective of synthetic degradation.

\section{Evaluation Strategy}
\label{sec:reg-evaluation}

\subsection{Main Metric: CI Overlap}
The primary measure of inferential utility is the Confidence Interval Overlap (CIO), defined as:
\begin{equation}
    \text{CIO} = \frac{1}{2} \left( \frac{\min(U_R, U_S) - \max(L_R, L_S)}{U_R - L_R} + \frac{\min(U_R, U_S) - \max(L_R, L_S)}{U_S - L_S} \right)
\end{equation}
where $[L_R, U_R]$ and $[L_S, U_S]$ are the 95\% CIs for the original and synthetic data, respectively.

\subsection{Supporting Metrics: Bias and CI Width}
To understand the \emph{source} of CIO degradation, we measure:
\begin{itemize}
    \item \textbf{Bias Ratio:} The ratio of the synthetic bias to the original standard error.
    \item \textbf{CI Width Ratio:} The ratio of the synthetic CI width to the original CI width.
\end{itemize}

\subsection{Univariate Analysis Strategy}
We first analyze each factor ($N, p, \rho, \dots$) independently by averaging the metrics across all other factors.

\subsection{Multivariate Analysis Strategy}
Finally, we use correlation heatmaps and multivariate plots to identify interactions between factors, such as how the impact of sample size changes as the number of predictors increases.
